% !TeX root = part3.tex

\documentclass[../final.tex]{subfiles}

\begin{document}

% При отдельной сборке продолжаем нумерацию полного конспекта.
\ifSubfilesClassLoaded{\setcounter{chapter}{2}\setcounter{section}{6}}{}


% ============================================================
% Лекция 15.09
% Численное интегрирование
% ============================================================

\chapter{Численное интегрирование}

\rsubtitle{Методы прямоугольников, трапеций, Симпсона и Рунге}

\captionsetup{
    font=footnotesize,
    labelfont=bf,
    skip=5pt
}

% Общий стиль рисунков лекции — тот же принцип, что и в part2.tex.
\tikzset{
    ovfaxis/.style={
        -{Stealth[length=2.2mm]},
        rtext,
        line width=0.85pt
    },
    ovfcurve/.style={
        rc4,
        line width=1.55pt
    },
    ovfapprox/.style={
        rtext!88,
        line width=1.0pt
    },
    ovfguide/.style={
        rtext!45,
        dashed,
        line width=0.70pt
    },
    ovfmesh/.style={
        rtext!55,
        line width=0.65pt
    },
    ovferror/.style={
        rc2,
        line width=0.95pt
    },
    ovfpoint/.style={
        circle,
        fill=rc3,
        draw=none,
        inner sep=1.65pt
    }
}

\pgfmathdeclarefunction{ovfintf}{1}{%
    \pgfmathparse{0.15 + 2.6*(1-exp(-0.52*(#1)))}%
}

% ============================================================
\section{Постановка задачи}
% ============================================================

Рассмотрим задачу численного вычисления определённого интеграла
%
\begin{equation*}
    I=\int_a^b f(x)\,dx.
\end{equation*}
%
Геометрически интеграл соответствует площади под графиком функции
$y=f(x)$ на отрезке $[a,b]$.

Разобьём интервал $[a,b]$ на $N$ одинаковых частей и введём
равномерную сетку
%
\begin{equation*}
    x_i=a+ih,
    \qquad
    i=0,1,\ldots,N,
\end{equation*}
%
где
%
\begin{equation*}
    h=\frac{b-a}{N}=\operatorname{const},
    \qquad
    x_0=a,
    \qquad
    x_N=b.
\end{equation*}

\begin{rbox}[left=18pt,right=18pt,top=14pt,bottom=10pt]{[]}
    \centering

    \begin{tikzpicture}[every node/.style={text=rtext},x=1.65cm,y=1.0cm]
        \path[pattern=north east lines,pattern color=rc3!65!rbg]
            (0.7,0) -- plot[domain=0.7:5.8,samples=100] (\x,{ovfintf(\x)})
            -- (5.8,0) -- cycle;
        \foreach \x in {0.7,1.55,2.4,3.25,4.1,4.95,5.8}{
            \draw[ovfguide] (\x,0) -- (\x,{ovfintf(\x)});
            \draw[ovfmesh] (\x,-0.055) -- (\x,0.055);
        }
        \draw[ovfaxis] (-0.15,0) -- (6.25,0) node[right] {$x$};
        \draw[ovfaxis] (0,-0.15) -- (0,3.35) node[left] {$y$};
        \draw[ovfcurve,domain=0.45:6.0,samples=100] plot (\x,{ovfintf(\x)});
        \foreach \x/\lab in {0.7/{a=x_0},1.55/{x_1},2.4/{x_2},3.25/{x_{i-1}},4.1/{x_i},4.95/{\cdots},5.8/{b=x_N}}
            \node[below=3pt] at (\x,0) {$\lab$};
        \draw[ovfmesh,{Stealth[length=1.5mm]}-{Stealth[length=1.5mm]}]
            (3.25,0.58) -- (4.1,0.58) node[midway,fill=rbg,inner sep=2pt] {$h$};
        \node[draw=rc3,circle,fill=rbg,text=rc3,inner sep=3pt] at (4.85,1.45) {$S$};
        \node[text=rc4,above=5pt] at (5.25,{ovfintf(5.25)}) {$y=f(x)$};
    \end{tikzpicture}

    \captionof{figure}{Разбиение отрезка интегрирования на равномерную сетку.}
    \label{fig:ovf_integral_grid}
\end{rbox}

Точное значение интеграла обозначим через $I$, а значение, полученное
численным методом на сетке с шагом $h$, --- через $S_h$. Тогда
%
\begin{equation*}
    I=S_h+R(h),
\end{equation*}
%
где
%
\begin{equation*}
    R(h)=I-S_h
\end{equation*}
%
--- остаточный член, или погрешность численного интегрирования.

\subsection{Квадратурные формулы как интерполяция}

На каждом маленьком интервале
%
\begin{equation*}
    [x_{i-1},x_i]
\end{equation*}
%
будем заменять исходную функцию $f(x)$ более простой функцией
$\varphi(x)$, которую легко проинтегрировать.

Если использовать полином степени $n=0$, то $\varphi(x)$ является
постоянной. Возможны три естественных варианта:
%
\begin{align*}
    \varphi(x)&=f(x_{i-1})
    &&\text{--- метод левых прямоугольников},\\[2pt]
    \varphi(x)&=f(x_i)
    &&\text{--- метод правых прямоугольников},\\[2pt]
    \varphi(x)&=f\!\left(\frac{x_{i-1}+x_i}{2}\right)
    &&\text{--- формула средних}.
\end{align*}

Если $n=1$, функцию заменяем линейным интерполяционным полиномом,
проходящим через значения на концах интервала:
%
\begin{equation*}
    \varphi(x)
    =
    \frac{x_i-x}{x_i-x_{i-1}}f(x_{i-1})
    +
    \frac{x-x_{i-1}}{x_i-x_{i-1}}f(x_i).
\end{equation*}
%
Интегрирование этой прямой приводит к \rconcept{методу трапеций}.

Если $n=2$, используем квадратичную интерполяцию по трём точкам.
В результате получаем \rconcept{формулу Симпсона}.


% ============================================================
\section{Методы прямоугольников}
% ============================================================

\subsection{Метод левых прямоугольников}

На каждом интервале $[x_{i-1},x_i]$ заменим функцию постоянным
значением в левой точке:
%
\begin{equation*}
    f(x)\approx f(x_{i-1}).
\end{equation*}
%
Тогда площадь под функцией на одном интервале приближается площадью
прямоугольника $f(x_{i-1})h$, а на всём отрезке
%
\begin{equation*}
    \rboxed{
    S_h^{(L)}
    =
    \sum_{i=1}^{N} f(x_{i-1})h
    }.
\end{equation*}

\begin{rbox}[left=18pt,right=18pt,top=14pt,bottom=10pt]{[]}
    \centering

    \begin{tikzpicture}[every node/.style={text=rtext,font=\small},x=0.82cm,y=1.0cm]
        % Общая сетка и увеличенная отдельная ячейка.
        \foreach \xa/\xb in {0.7/1.55,1.55/2.4,2.4/3.25,3.25/4.1,4.1/4.95,4.95/5.8}{
            \pgfmathsetmacro{\yy}{ovfintf(\xa)}
            \path[pattern=north east lines,pattern color=rc3!65!rbg] (\xa,0) rectangle (\xb,\yy);
            \draw[ovfapprox] (\xa,0) -- (\xa,\yy) -- (\xb,\yy) -- (\xb,0);
        }
        \draw[ovfaxis] (-0.15,0) -- (6.35,0) node[right] {$x$};
        \draw[ovfaxis] (0,-0.15) -- (0,3.2) node[left] {$y$};
        \draw[ovfcurve,domain=0.45:6.0,samples=100] plot (\x,{ovfintf(\x)});
        \node[below=3pt] at (0.7,0) {$a$};
        \node[below=3pt] at (5.8,0) {$b$};
        \node[text=rc4,above=4pt] at (4.5,{ovfintf(4.5)}) {$y=f(x)$};
        \begin{scope}[shift={(9.4,0)},x=1.13cm]
            \pgfmathsetmacro{\ya}{ovfintf(1)}
            \pgfmathsetmacro{\yb}{ovfintf(3.8)}
            \pgfmathsetmacro{\yr}{ovfintf(1.0)}
            \path[pattern=north east lines,pattern color=rc3!65!rbg] (1,0) rectangle (3.8,\yr);
            \path[fill=rc2,fill opacity=0.23] (1,\yr)
                -- plot[domain=1:3.8,samples=80] (\x,{ovfintf(\x)}) -- (3.8,\yr) -- cycle;
            \draw[ovfguide] (0,\ya) -- (1,\ya);
            \draw[ovfguide] (0,\yb) -- (3.8,\yb);
            \draw[ovfguide] (1,0) -- (1,\ya);
            \draw[ovfguide] (3.8,0) -- (3.8,\yb);
            \draw[ovfapprox] (1,0) -- (1,\yr) -- (3.8,\yr) -- (3.8,0);
            \draw[ovfaxis] (-0.15,0) -- (4.55,0) node[right] {$x$};
            \draw[ovfaxis] (0,-0.15) -- (0,3.2) node[left] {$y$};
            \draw[ovfcurve,domain=0.55:4.2,samples=80] plot (\x,{ovfintf(\x)});
            \node[ovfpoint] at (1.0,\yr) {};
            \node[left=3pt] at (0,\ya) {$f(x_{i-1})$};
            \node[left=3pt] at (0,\yb) {$f(x_i)$};
            \node[below=3pt] at (1,0) {$x_{i-1}$};
            \node[below=3pt] at (3.8,0) {$x_i$};
            \draw[ovfmesh,{Stealth[length=1.4mm]}-{Stealth[length=1.4mm]}]
                (1,0.45) -- (3.8,0.45) node[midway,fill=rbg,inner sep=2pt] {$h$};
            \node[text=rc2,anchor=west] (err) at (2.8,3.05) {$R_i^{(L)}>0$};
            \draw[ovferror,-{Stealth[length=1.4mm]}] (err.south) to[bend right=15] (3.0,1.65);
            \node[fill=rbg,inner sep=2pt,text=rc3] at (2.4,0.95) {$S_i$};
        \end{scope}
    \end{tikzpicture}

    \captionof{figure}{Метод левых прямоугольников.}
    \label{fig:ovf_left_rectangles}
\end{rbox}

Точное значение интеграла представим как сумму интегралов по ячейкам:
%
\begin{equation*}
    I
    =
    \int_a^b f(x)\,dx
    =
    \sum_{i=1}^{N}
    \int_{x_{i-1}}^{x_i} f(x)\,dx.
\end{equation*}
%
Поэтому
%
\begin{align*}
    R_L(h)
    &=I-S_h^{(L)}\\
    &=
    \sum_{i=1}^{N}
    \left[
        \int_{x_{i-1}}^{x_i} f(x)\,dx
        -f(x_{i-1})h
    \right].
\end{align*}

Чтобы оценить погрешность, разложим функцию в ряд Тейлора около
левой точки $x_{i-1}$:
%
\begin{equation*}
    f(x)
    =
    f(x_{i-1})
    +f'(x_{i-1})(x-x_{i-1})
    +\frac12 f''(x_{i-1})(x-x_{i-1})^2
    +O(h^3).
\end{equation*}

Интегрируя по одному интервалу, получаем
%
\begin{align*}
    \int_{x_{i-1}}^{x_i}f(x)\,dx
    &=
    f(x_{i-1})h
    +f'(x_{i-1})
    \int_{x_{i-1}}^{x_i}(x-x_{i-1})\,dx
    +O(h^3)\\
    &=
    f(x_{i-1})h
    +\frac12 f'(x_{i-1})h^2
    +O(h^3).
\end{align*}

После вычитания площади прямоугольника остаётся локальная ошибка
%
\begin{equation*}
    R_i
    =
    \frac12 f'(x_{i-1})h^2+O(h^3).
\end{equation*}

Суммируя по всей сетке,
%
\begin{align*}
    R_L(h)
    &=
    \sum_{i=1}^{N}
    \left[
        \frac12 f'(x_{i-1})h^2+O(h^3)
    \right]\\
    &=
    \frac{h}{2}
    \sum_{i=1}^{N}f'(x_{i-1})h
    +O(h^2).
\end{align*}

При $h\to0$ сумма переходит в интеграл:
%
\begin{equation*}
    \sum_{i=1}^{N}f'(x_{i-1})h
    \longrightarrow
    \int_a^b f'(x)\,dx.
\end{equation*}

Следовательно,
%
\begin{equation*}
    \rboxed{
    R_L(h)
    =
    \frac{h}{2}\int_a^b f'(x)\,dx
    +O(h^2)
    }
    =O(h).
\end{equation*}

Таким образом, метод левых прямоугольников имеет
\rresult{первый порядок точности}.


\subsection{Метод правых прямоугольников}

Теперь на каждом интервале используем значение функции в правом конце:
%
\begin{equation*}
    f(x)\approx f(x_i).
\end{equation*}
%
Получаем
%
\begin{equation*}
    \rboxed{
    S_h^{(R)}
    =
    \sum_{i=1}^{N} f(x_i)h
    }.
\end{equation*}

\begin{rbox}[left=18pt,right=18pt,top=14pt,bottom=10pt]{[]}
    \centering

    \begin{tikzpicture}[every node/.style={text=rtext,font=\small},x=0.82cm,y=1.0cm]
        % Общая сетка и увеличенная отдельная ячейка.
        \foreach \xa/\xb in {0.7/1.55,1.55/2.4,2.4/3.25,3.25/4.1,4.1/4.95,4.95/5.8}{
            \pgfmathsetmacro{\yy}{ovfintf(\xb)}
            \path[pattern=north east lines,pattern color=rc3!65!rbg] (\xa,0) rectangle (\xb,\yy);
            \draw[ovfapprox] (\xa,0) -- (\xa,\yy) -- (\xb,\yy) -- (\xb,0);
        }
        \draw[ovfaxis] (-0.15,0) -- (6.35,0) node[right] {$x$};
        \draw[ovfaxis] (0,-0.15) -- (0,3.2) node[left] {$y$};
        \draw[ovfcurve,domain=0.45:6.0,samples=100] plot (\x,{ovfintf(\x)});
        \node[below=3pt] at (0.7,0) {$a$};
        \node[below=3pt] at (5.8,0) {$b$};
        \node[text=rc4,above=4pt] at (4.5,{ovfintf(4.5)}) {$y=f(x)$};
        \begin{scope}[shift={(9.4,0)},x=1.13cm]
            \pgfmathsetmacro{\ya}{ovfintf(1)}
            \pgfmathsetmacro{\yb}{ovfintf(3.8)}
            \pgfmathsetmacro{\yr}{ovfintf(3.8)}
            \path[pattern=north east lines,pattern color=rc3!65!rbg] (1,0) rectangle (3.8,\yr);
            \path[fill=rc2,fill opacity=0.23] (1,\yr)
                -- plot[domain=1:3.8,samples=80] (\x,{ovfintf(\x)}) -- (3.8,\yr) -- cycle;
            \draw[ovfguide] (0,\ya) -- (1,\ya);
            \draw[ovfguide] (0,\yb) -- (3.8,\yb);
            \draw[ovfguide] (1,0) -- (1,\ya);
            \draw[ovfguide] (3.8,0) -- (3.8,\yb);
            \draw[ovfapprox] (1,0) -- (1,\yr) -- (3.8,\yr) -- (3.8,0);
            \draw[ovfaxis] (-0.15,0) -- (4.55,0) node[right] {$x$};
            \draw[ovfaxis] (0,-0.15) -- (0,3.2) node[left] {$y$};
            \draw[ovfcurve,domain=0.55:4.2,samples=80] plot (\x,{ovfintf(\x)});
            \node[ovfpoint] at (3.8,\yr) {};
            \node[left=3pt] at (0,\ya) {$f(x_{i-1})$};
            \node[left=3pt] at (0,\yb) {$f(x_i)$};
            \node[below=3pt] at (1,0) {$x_{i-1}$};
            \node[below=3pt] at (3.8,0) {$x_i$};
            \draw[ovfmesh,{Stealth[length=1.4mm]}-{Stealth[length=1.4mm]}]
                (1,0.45) -- (3.8,0.45) node[midway,fill=rbg,inner sep=2pt] {$h$};
            \node[text=rc2,anchor=west] (err) at (2.8,3.05) {$R_i^{(R)}<0$};
            \draw[ovferror,-{Stealth[length=1.4mm]}] (err.south) to[bend right=15] (1.6,1.95);
            \node[fill=rbg,inner sep=2pt,text=rc3] at (2.4,0.95) {$S_i$};
        \end{scope}
    \end{tikzpicture}

    \captionof{figure}{Метод правых прямоугольников.}
    \label{fig:ovf_right_rectangles}
\end{rbox}

Разложение относительно правой точки $x_i$ приводит к тому же по
модулю, но противоположному по знаку главному члену погрешности:
%
\begin{equation*}
    \rboxed{
    R_R(h)
    =
    -\frac{h}{2}\int_a^b f'(x)\,dx
    +O(h^2)
    }
    =O(h).
\end{equation*}

Следовательно, метод правых прямоугольников также имеет первый
порядок точности. При этом
%
\begin{equation*}
    R_L(h)
    \simeq
    +\frac{h}{2}\int_a^b f'(x)\,dx,
    \qquad
    R_R(h)
    \simeq
    -\frac{h}{2}\int_a^b f'(x)\,dx.
\end{equation*}
%
Главные части погрешностей имеют одинаковый модуль и противоположные
знаки. Это позволяет построить метод более высокого порядка точности.


% ============================================================
\section{Метод трапеций}
% ============================================================

Для левых и правых прямоугольников имеем
%
\begin{equation*}
    I=S_L+R_L(h),
    \qquad
    I=S_R+R_R(h).
\end{equation*}

Сложим эти равенства и разделим результат на два:
%
\begin{equation*}
    I
    =
    \frac{S_L+S_R}{2}
    +
    \frac{R_L(h)+R_R(h)}{2}.
\end{equation*}

Поскольку главные члены погрешностей порядка $O(h)$ имеют
противоположные знаки, они сокращаются. В результате
%
\begin{equation*}
    I
    =
    \sum_{i=1}^{N}
    \frac{f(x_{i-1})+f(x_i)}{2}h
    +O(h^2).
\end{equation*}

Таким образом, получаем квадратурную формулу трапеций
%
\begin{equation*}
    \rboxed{
    S_h^{(T)}
    =
    \sum_{i=1}^{N}
    \frac{f(x_{i-1})+f(x_i)}{2}\,h
    }.
\end{equation*}

\begin{rbox}[left=18pt,right=18pt,top=14pt,bottom=10pt]{[]}
    \centering

    \begin{tikzpicture}[every node/.style={text=rtext},x=1.55cm,y=1.1cm]
        \pgfmathsetmacro{\ya}{ovfintf(1)}
        \pgfmathsetmacro{\yb}{ovfintf(3.8)}
        \path[pattern=north east lines,pattern color=rc3!65!rbg]
            (1,0) -- (1,\ya) -- (3.8,\yb) -- (3.8,0) -- cycle;
        \draw[ovfguide] (0,\ya) -- (1,\ya);
        \draw[ovfguide] (0,\yb) -- (3.8,\yb);
        \draw[ovfapprox] (1,0) -- (1,\ya) -- (3.8,\yb) -- (3.8,0);
        \draw[ovfaxis] (-0.1,0) -- (4.65,0) node[right] {$x$};
        \draw[ovfaxis] (0,-0.1) -- (0,3.15) node[left] {$y$};
        \draw[ovfcurve,domain=0.55:4.25,samples=100] plot (\x,{ovfintf(\x)});
        \node[ovfpoint] at (1,\ya) {};
        \node[ovfpoint] at (3.8,\yb) {};
        \node[left=4pt] at (0,\ya) {$f(x_{i-1})$};
        \node[left=4pt] at (0,\yb) {$f(x_i)$};
        \node[below=3pt] at (1,0) {$x_{i-1}$};
        \node[below=3pt] at (3.8,0) {$x_i$};
        \node[text=rc4,above=5pt] at (3.85,{ovfintf(3.85)}) {$y=f(x)$};
        \node[fill=rbg,inner sep=3pt,text=rc3] (trap) at (2.45,0.85) {трапеция};
        \draw[rtext!85,-{Stealth[length=1.6mm]}] (trap.north) -- (2.9,1.98);
        \draw[ovfmesh,{Stealth[length=1.5mm]}-{Stealth[length=1.5mm]}]
            (1,0.3) -- (3.8,0.3) node[midway,fill=rbg,inner sep=2pt] {$h$};
    \end{tikzpicture}

    \captionof{figure}{На каждом интервале график функции заменяется хордой.}
    \label{fig:ovf_trapezoid}
\end{rbox}

\subsection{Погрешность метода трапеций}

Обозначим середину $i$-го интервала через
%
\begin{equation*}
    c_i=\frac{x_{i-1}+x_i}{2}.
\end{equation*}
%
Тогда
%
\begin{equation*}
    x_{i-1}=c_i-\frac h2,
    \qquad
    x_i=c_i+\frac h2.
\end{equation*}

Разложим функцию около середины интервала:
%
\begin{equation*}
    f(x)
    =
    f(c_i)
    +f'(c_i)(x-c_i)
    +\frac12 f''(c_i)(x-c_i)^2
    +O(h^3).
\end{equation*}

При интегрировании по симметричному промежутку линейный член исчезает:
%
\begin{equation*}
    \int_{-h/2}^{h/2}t\,dt=0.
\end{equation*}
%
Поэтому
%
\begin{align*}
    \int_{x_{i-1}}^{x_i}f(x)\,dx
    &=
    f(c_i)h
    +\frac12 f''(c_i)
    \int_{-h/2}^{h/2}t^2\,dt
    +O(h^5)\\
    &=
    f(c_i)h
    +\frac{1}{24}f''(c_i)h^3
    +O(h^5).
\end{align*}

Теперь разложим значения на концах интервала:
%
\begin{align*}
    f\!\left(c_i-\frac h2\right)
    &=
    f(c_i)-\frac h2 f'(c_i)
    +\frac{h^2}{8}f''(c_i)
    +O(h^3),\\
    f\!\left(c_i+\frac h2\right)
    &=
    f(c_i)+\frac h2 f'(c_i)
    +\frac{h^2}{8}f''(c_i)
    +O(h^3).
\end{align*}

Следовательно,
%
\begin{equation*}
    \frac h2
    \left[
        f(x_{i-1})+f(x_i)
    \right]
    =
    f(c_i)h
    +\frac18 f''(c_i)h^3
    +O(h^5).
\end{equation*}

Локальная погрешность равна
%
\begin{equation*}
    R_i^{(T)}
    =
    -\frac1{12}f''(c_i)h^3
    +O(h^5).
\end{equation*}

После суммирования по всей сетке
%
\begin{align*}
    R_T(h)
    &=
    -\frac{h^2}{12}
    \sum_{i=1}^{N}f''(c_i)h
    +O(h^4)\\
    &=
    \rboxed{
    -\frac{h^2}{12}\int_a^b f''(x)\,dx
    +O(h^4)
    }.
\end{align*}

Следовательно,
%
\begin{equation*}
    R_T(h)=O(h^2),
\end{equation*}
%
и метод трапеций имеет \rresult{второй порядок точности}.


% ============================================================
\section{Формула средних}
% ============================================================

В методе средних функцию на каждом интервале заменяем её значением
в центральной точке
%
\begin{equation*}
    c_i=\frac{x_{i-1}+x_i}{2}.
\end{equation*}

Получаем
%
\begin{equation*}
    \rboxed{
    S_h^{(M)}
    =
    \sum_{i=1}^{N}
    f\!\left(\frac{x_{i-1}+x_i}{2}\right)h
    }.
\end{equation*}

\begin{rbox}[left=18pt,right=18pt,top=14pt,bottom=10pt]{[]}
    \centering

    \begin{tikzpicture}[every node/.style={text=rtext},x=1.55cm,y=1.1cm]
        \pgfmathsetmacro{\ym}{ovfintf(2.4)}
        \path[pattern=north east lines,pattern color=rc3!65!rbg] (1,0) rectangle (3.8,\ym);
        \path[fill=rc2,fill opacity=0.28] (1,\ym)
            -- plot[domain=1:2.4,samples=50] (\x,{ovfintf(\x)}) -- cycle;
        \path[fill=rc2,fill opacity=0.28] (2.4,\ym)
            -- plot[domain=2.4:3.8,samples=50] (\x,{ovfintf(\x)}) -- (3.8,\ym) -- cycle;
        \draw[ovfguide] (0,\ym) -- (1,\ym);
        \draw[ovfguide] (2.4,0) -- (2.4,\ym);
        \draw[ovfguide] (3.8,\ym) -- (3.8,{ovfintf(3.8)});
        \draw[ovfapprox] (1,0) -- (1,\ym) -- (3.8,\ym) -- (3.8,0);
        \draw[ovfaxis] (-0.1,0) -- (4.65,0) node[right] {$x$};
        \draw[ovfaxis] (0,-0.1) -- (0,3.15) node[left] {$y$};
        \draw[ovfcurve,domain=0.55:4.25,samples=100] plot (\x,{ovfintf(\x)});
        \node[ovfpoint] at (2.4,\ym) {};
        \node[left=4pt] at (0,\ym) {$f(c_i)$};
        \node[below=3pt] at (1,0) {$x_{i-1}$};
        \node[below=3pt] at (3.8,0) {$x_i$};
        \node[below=3pt] at (2.4,0) {$c_i$};
        \node[text=rc4,above=5pt] at (3.9,{ovfintf(3.9)}) {$y=f(x)$};
        \node[text=rc2] (minus) at (0.95,2.6) {$-$};
        \draw[ovferror,-{Stealth[length=1.4mm]}] (minus.south) -- (1.45,1.77);
        \node[text=rc2] (plus) at (3.1,2.9) {$+$};
        \draw[ovferror,-{Stealth[length=1.4mm]}] (plus.south) -- (3.25,2.11);
        \draw[ovfmesh,{Stealth[length=1.4mm]}-{Stealth[length=1.4mm]}]
            (1,0.52) -- (2.4,0.52) node[midway,fill=rbg,inner sep=2pt] {$h/2$};
        \draw[ovfmesh,{Stealth[length=1.4mm]}-{Stealth[length=1.4mm]}]
            (2.4,0.52) -- (3.8,0.52) node[midway,fill=rbg,inner sep=2pt] {$h/2$};
    \end{tikzpicture}

    \captionof{figure}{Формула средних: высота прямоугольника берётся в центре интервала.}
    \label{fig:ovf_midpoint}
\end{rbox}

Остаточный член имеет вид
%
\begin{equation*}
    R_M(h)
    =
    \sum_{i=1}^{N}
    \left[
        \int_{x_{i-1}}^{x_i}f(x)\,dx
        -f(c_i)h
    \right].
\end{equation*}

Снова используем разложение около середины:
%
\begin{equation*}
    f(x)
    =
    f(c_i)
    +f'(c_i)(x-c_i)
    +\frac12 f''(c_i)(x-c_i)^2
    +O(h^3).
\end{equation*}

Нечётный член исчезает при интегрировании, поэтому
%
\begin{align*}
    \int_{x_{i-1}}^{x_i}f(x)\,dx-f(c_i)h
    &=
    \frac12 f''(c_i)
    \int_{-h/2}^{h/2}t^2\,dt
    +O(h^5)\\
    &=
    \frac1{24}f''(c_i)h^3
    +O(h^5).
\end{align*}

Суммируя по всем интервалам,
%
\begin{align*}
    R_M(h)
    &=
    \frac{h^2}{24}
    \sum_{i=1}^{N}f''(c_i)h
    +O(h^4)\\
    &=
    \rboxed{
    \frac{h^2}{24}\int_a^b f''(x)\,dx
    +O(h^4)
    }.
\end{align*}

Следовательно,
%
\begin{equation*}
    R_M(h)=O(h^2),
\end{equation*}
%
то есть формула средних имеет второй порядок точности.

Сравним главные члены погрешностей формулы трапеций и формулы средних:
%
\begin{equation*}
    R_T(h)
    \simeq
    -\frac{h^2}{12}\int_a^b f''(x)\,dx,
    \qquad
    R_M(h)
    \simeq
    +\frac{h^2}{24}\int_a^b f''(x)\,dx.
\end{equation*}
%
Их знаки противоположны, причём
%
\begin{equation*}
    R_T(h)\simeq -2R_M(h).
\end{equation*}


% ============================================================
\section{Метод Симпсона}
% ============================================================

Используем найденное соотношение между погрешностями трапеций и
средних. Имеем
%
\begin{equation*}
    I=S_T+R_T,
    \qquad
    I=S_M+R_M.
\end{equation*}

Умножим второе равенство на $2$ и сложим с первым:
%
\begin{equation*}
    3I
    =
    S_T+2S_M+R_T+2R_M.
\end{equation*}

Но главные члены второго порядка взаимно уничтожаются:
%
\begin{align*}
    R_T+2R_M
    &\simeq
    -\frac{h^2}{12}\int_a^b f''(x)\,dx
    +2\cdot\frac{h^2}{24}\int_a^b f''(x)\,dx\\
    &=0.
\end{align*}

Следовательно,
%
\begin{equation*}
    I
    =
    \frac{S_T+2S_M}{3}
    +O(h^4).
\end{equation*}

Подставляя формулы трапеций и средних, получаем
%
\begin{align*}
    S_h^{(S)}
    &=
    \frac13
    \sum_{i=1}^{N}
    \left[
        2f\!\left(\frac{x_{i-1}+x_i}{2}\right)
        +\frac{f(x_{i-1})+f(x_i)}{2}
    \right]h\\[3pt]
    &=
    \rboxed{
    \sum_{i=1}^{N}
    \frac{
        f(x_{i-1})
        +4f\!\left(\dfrac{x_{i-1}+x_i}{2}\right)
        +f(x_i)
    }{6}\,h
    }.
\end{align*}

\begin{rbox}[left=18pt,right=18pt,top=14pt,bottom=10pt]{[]}
    \centering

    \begin{tikzpicture}[every node/.style={text=rtext},x=1.55cm,y=1.12cm]
        \def\xa{0.9}
        \def\xb{3.7}
        \def\xm{2.3}
        \pgfmathsetmacro{\ya}{ovfintf(\xa)}
        \pgfmathsetmacro{\ym}{ovfintf(\xm)}
        \pgfmathsetmacro{\yb}{ovfintf(\xb)}
        \foreach \xx/\yy/\lab in {\xa/\ya/{f(x_{i-1})},\xm/\ym/{f(c_i)},\xb/\yb/{f(x_i)}}{
            \draw[ovfguide] (0,\yy) -- (\xx,\yy) -- (\xx,0);
            \node[left=4pt,font=\small] at (0,\yy) {$\lab$};
        }
        \draw[ovfaxis] (-0.1,0) -- (4.55,0) node[right] {$x$};
        \draw[ovfaxis] (0,-0.1) -- (0,3.15) node[left] {$y$};
        \draw[ovfcurve,domain=0.5:4.15,samples=100] plot (\x,{ovfintf(\x)});
        \draw[ovfapprox,densely dashed,domain=\xa:\xb,samples=80]
            plot (\x,{\ya*((\x-\xm)*(\x-\xb))/((\xa-\xm)*(\xa-\xb))
                +\ym*((\x-\xa)*(\x-\xb))/((\xm-\xa)*(\xm-\xb))
                +\yb*((\x-\xa)*(\x-\xm))/((\xb-\xa)*(\xb-\xm))});
        \foreach \xx/\yy in {\xa/\ya,\xm/\ym,\xb/\yb}
            \node[ovfpoint] at (\xx,\yy) {};
        \node[below=3pt] at (\xa,0) {$x_{i-1}$};
        \node[below=3pt] at (\xm,0) {$c_i$};
        \node[below=3pt] at (\xb,0) {$x_i$};
        \node[text=rc4,above=5pt] at (3.75,{ovfintf(3.75)}) {$y=f(x)$};
        \node (quad) at (3.45,1.12) {$\varphi_2(x)$};
        \draw[rtext!85,-{Stealth[length=1.4mm]}] (quad.north) -- (3.1,2.23);
        \draw[ovfmesh,{Stealth[length=1.4mm]}-{Stealth[length=1.4mm]}]
            (\xa,0.38) -- (\xb,0.38) node[midway,fill=rbg,inner sep=2pt] {$h$};
    \end{tikzpicture}

    \captionof{figure}{В формуле Симпсона функция заменяется квадратичной интерполяцией по трём точкам.}
    \label{fig:ovf_simpson}
\end{rbox}

Для остаточного члена формулы Симпсона получаем
%
\begin{equation*}
    \rboxed{
    R_S(h)
    =
    -\frac{h^4}{2880}
    \int_a^b f^{(4)}(x)\,dx
    +O(h^6)
    }.
\end{equation*}

Следовательно,
%
\begin{equation*}
    R_S(h)=O(h^4),
\end{equation*}
%
и метод Симпсона имеет \rresult{четвёртый порядок точности}.

\remark
В этой записи $h=x_i-x_{i-1}$ --- ширина целого интервала, на котором
используются две крайние точки и его середина. Поэтому коэффициент в
остаточном члене записан как $1/2880$.


% ============================================================
\section{Метод Рунге}
% ============================================================

Метод Рунге позволяет повысить порядок точности, используя результаты
расчёта одной и той же величины на двух сетках с разными шагами.

Пусть численный метод имеет порядок точности $p$. Тогда результат на
сетке с шагом $h$ можно записать в виде
%
\begin{equation}
    I
    =
    S(h)
    +C h^p
    +O(h^{p+1}),
    \label{eq:ovf_runge_h}
\end{equation}
%
где $C=\operatorname{const}$ относительно $h$.

Теперь изменим шаг сетки в $r$ раз:
%
\begin{equation*}
    h\longrightarrow rh.
\end{equation*}
%
Для новой сетки имеем
%
\begin{equation}
    I
    =
    S(rh)
    +C(rh)^p
    +O(h^{p+1})
    =
    S(rh)
    +Cr^ph^p
    +O(h^{p+1}).
    \label{eq:ovf_runge_rh}
\end{equation}

Умножим \eqref{eq:ovf_runge_h} на $r^p$:
%
\begin{equation*}
    r^pI
    =
    r^pS(h)
    +Cr^ph^p
    +O(h^{p+1}).
\end{equation*}

Теперь вычтем \eqref{eq:ovf_runge_rh}. Главный член погрешности
$Cr^ph^p$ сокращается:
%
\begin{align*}
    r^pI-I
    &=
    r^pS(h)-S(rh)
    +O(h^{p+1}),\\
    (r^p-1)I
    &=
    r^pS(h)-S(rh)
    +O(h^{p+1}).
\end{align*}

Отсюда получаем формулу Рунге
%
\begin{equation*}
    \rboxed{
    I
    =
    \frac{r^pS(h)-S(rh)}{r^p-1}
    +O(h^{p+1})
    }.
\end{equation*}

Таким образом, линейная комбинация результатов на двух сетках позволяет
убрать главный член погрешности $Ch^p$.

Если обозначить уточнённое значение через $S_R(h)$, то
%
\begin{equation*}
    \rboxed{
    S_R(h)
    =
    \frac{r^pS(h)-S(rh)}{r^p-1}
    }.
\end{equation*}

\begin{rbox}[left=18pt,right=18pt,top=13pt,bottom=13pt]{<u}
    \textbf{Идея метода Рунге.}
    Если известна структура главного члена погрешности
    $R(h)=Ch^p+O(h^{p+1})$, то два расчёта с шагами $h$ и $rh$
    позволяют исключить неизвестную константу $C$ и получить более
    точное приближение к интегралу.
\end{rbox}

\end{document}
